\chapter{Análisis del problema}
\label{cap:analisis_problema}
En este capítulo se presenta la problemática que da origen al Trabajo Terminal. En primer lugar, se describe el problema general: la ineficiencia del procedimiento de atención de quejas de la Defensoría de los Derechos Politécnicos. Posteriormente, se desglosa en los problemas específicos que lo conforman y se analizan las causas que los originan. Finalmente, se presenta la propuesta de solución que responde a cada problema específico; los objetivos, el detalle de la propuesta y los riesgos y supuestos del proyecto se presentan en el Capítulo~\ref{cap:proyecto}.

\section{Dominio}
\label{sec:dominio}
El Trabajo Terminal se ubica en el ámbito de la protección de los derechos de la comunidad del Instituto Politécnico Nacional (IPN), institución pública de educación media superior y superior cuya organización, funcionamiento y atribuciones se encuentran regulados por su marco normativo \cite{ipn_ley_organica}. Dentro de su estructura, la Defensoría de los Derechos Politécnicos (DDP) fue creada mediante el acuerdo publicado en la Gaceta Politécnica número 622, en enero de 2006, como el órgano encargado de promover, proteger y defender los derechos de los integrantes de la comunidad politécnica \cite{ipn_acuerdo622}. Las funciones y responsabilidades de sus áreas se establecen en su Manual de Organización \cite{ipn_manual_organizacion_ddp}, mientras que la forma en que debe atenderse cada asunto se describe en su Manual de Procedimientos \cite{ipn_manual_procedimientos_ddp}.

	La DDP atiende a estudiantes, personal académico, personal de apoyo y demás integrantes de la comunidad politécnica. El instrumento central de su labor es la \emph{queja}: un escrito mediante el cual una persona, denominada \emph{quejoso}, hace del conocimiento de la Defensoría hechos que considera una posible vulneración de sus derechos, atribuidos a una persona o a un área del Instituto, a quien se denomina \emph{persona señalada}. Cada queja admitida da origen a un \emph{expediente}, que concentra la narrativa de los hechos, los datos del quejoso, las evidencias presentadas y los documentos que se generan durante su atención, como oficios, citas, notas de análisis y determinaciones. El expediente es, por tanto, el objeto de información sobre el que trabajan todas las áreas de la Defensoría.

	La atención de una queja sigue un procedimiento institucional en el que intervienen distintas áreas de manera sucesiva \cite{ipn_manual_procedimientos_ddp}:

	\begin{itemize}
		\item \textbf{Recepción.} Recibe la queja, verifica que cumpla los requisitos de admisión y, cuando es procedente, la registra y la turna al área correspondiente; en caso contrario, la rechaza o solicita al quejoso que la corrija o complemente.
		\item \textbf{Primer Contacto.} Analiza los hechos descritos, consulta los antecedentes relacionados, puede citar al quejoso o solicitar la versión de la persona señalada y determina si el asunto es competencia de la Defensoría; cuando no lo es, orienta al quejoso o remite el asunto a la instancia competente.
		\item \textbf{Subdefensoría.} Investiga los asuntos admitidos y da seguimiento con las unidades involucradas mediante solicitudes de información sujetas a plazos: la primera solicitud a una unidad cuenta con un plazo de diez días hábiles y las solicitudes subsecuentes, de cinco días hábiles.
		\item \textbf{Titular de la Defensoría.} Supervisa el conjunto de los casos, participa en las etapas de decisión y resolución de los expedientes y coordina la operación de las áreas.
	\end{itemize}

	El dominio presenta características que condicionan cualquier solución tecnológica. En primer lugar, el procedimiento es secuencial: cada área depende de la información generada por la anterior, por lo que la disponibilidad y la integridad del expediente determinan la continuidad de la atención. En segundo lugar, la información principal de cada caso se expresa en lenguaje natural, ya que las narrativas son redactadas libremente por los quejosos. En tercer lugar, los expedientes contienen datos personales y descripciones de situaciones sensibles, por lo que su consulta debe restringirse al personal autorizado según sus responsabilidades. Finalmente, las determinaciones sobre la competencia, la investigación y la resolución de cada asunto corresponden al criterio jurídico del personal de la Defensoría y no pueden delegarse en un sistema informático; cualquier herramienta de apoyo debe, por tanto, auxiliar al personal sin sustituir sus decisiones.

	\section{Problema general}

	El problema central es que este procedimiento resulta ineficiente para el volumen de quejas que recibe la Defensoría. Su ejecución recae principalmente en actividades manuales y secuenciales: el registro y la consulta de los expedientes, la lectura completa de cada narrativa, la búsqueda de antecedentes y la comunicación del avance entre áreas dependen del tiempo y la intervención directa del personal. En consecuencia, el tiempo de atención de cada caso se incrementa, las quejas se acumulan en las distintas etapas y la información necesaria para darles continuidad no siempre está disponible de manera oportuna para quien la requiere. Las dificultades particulares que conforman esta problemática se describen en el siguiente apartado.

	\section{Problemas específicos}
	A partir del levantamiento de requerimientos realizado con el personal de la Defensoría y del análisis de los procedimientos institucionales, se identificaron los siguientes problemas, que en conjunto explican la ineficiencia del proceso descrito.

	\begin{enumerate}
		\item \label{pe:volumen}\textbf{Volumen de quejas superior a la capacidad de atención:}
		De acuerdo con lo señalado por el personal de la Defensoría, el número de quejas que ingresan supera la capacidad del personal para atenderlas y concluirlas en los plazos institucionales. Como cada etapa del procedimiento depende de actividades manuales, el tiempo invertido en tareas de registro, consulta y actualización reduce el tiempo disponible para el análisis de los casos, lo que provoca la acumulación de expedientes pendientes y el retraso en la respuesta a los quejosos.

		\item \label{pe:narrativas}\textbf{Revisión de narrativas extensas:}
		Las quejas contienen descripciones redactadas libremente por los quejosos, que con frecuencia son extensas, repetitivas o no siguen un orden cronológico. Para continuar con el procedimiento, el personal debe leer cada narrativa completa e identificar por sí mismo los hechos relevantes, lo que consume una parte importante del tiempo de revisión de cada expediente.

		\item \label{pe:antecedentes}\textbf{Búsqueda manual de antecedentes:}
		Durante la revisión de una queja es necesario saber si existen expedientes anteriores relacionados con los mismos hechos, personas o unidades. Esta búsqueda se realiza de forma manual y depende de la memoria del personal o de la revisión de registros previos, lo que resulta lento y poco confiable cuando los casos similares están descritos con palabras distintas.

		\item \label{pe:duplicadas}\textbf{Duplicidad de quejas:}
		De acuerdo con lo señalado por el personal de la Defensoría, existen quejosos que presentan nuevamente una solicitud sobre hechos ya reportados, particularmente cuando consideran que no han recibido respuesta. Al no contar con un mecanismo para identificarlas, estas quejas pueden registrarse y atenderse como casos nuevos, lo que duplica el trabajo y fragmenta la información de un mismo asunto.

		\item \label{pe:estado}\textbf{Falta de seguimiento de la queja:}
		No existe un medio único para conocer en qué etapa se encuentra una queja, qué actuaciones se han realizado y qué plazos están por vencer. Esto dificulta priorizar los casos urgentes y obliga a los quejosos a solicitar información sobre el avance de su queja.

		\item \label{pe:informacion}\textbf{Falta de coordinación entre áreas:}
		Cuando un expediente pasa de un área a otra, la información generada en la etapa anterior ---documentos, oficios, citas y determinaciones--- no siempre se encuentra disponible e integrada para quien continúa el procedimiento. Esto provoca consultas adicionales, retrasos y dificultades de coordinación entre el personal.

	\end{enumerate}

	Estos problemas se manifiestan en distintas etapas del procedimiento, pero mantienen una relación directa con la organización de la información, la continuidad de los expedientes y la carga operativa del personal. Su análisis permite identificar los factores que contribuyen a su aparición, los cuales se desarrollan en el siguiente apartado.

	\section{Causas}
	\label{sec:causas}
	Las problemáticas identificadas en la gestión de quejas de la Defensoría de los Derechos Politécnicos se relacionan con diversos factores administrativos, operativos y tecnológicos que intervienen en el desarrollo de sus procedimientos internos. A partir del análisis de los procesos institucionales y de la información obtenida durante el levantamiento de requerimientos, se reconocen condiciones que pueden contribuir a las dificultades de organización, consulta y seguimiento de los expedientes. Las causas se presentan en el mismo orden que los problemas específicos, indicando en cada una el problema que explica.

	\textbf{Dependencia de actividades manuales (problema~\ref{pe:volumen}).} La dependencia de actividades manuales constituye uno de los factores que contribuyen a las dificultades operativas durante la atención de quejas. Los procedimientos institucionales requieren que el personal realice tareas de registro, revisión, actualización y consulta de información en distintas etapas del proceso \cite{ipn_manual_procedimientos_ddp}. Cuando estas actividades dependen principalmente de la intervención humana y carecen de mecanismos que faciliten su ejecución, el tiempo disponible del personal debe distribuirse entre tareas administrativas y actividades que requieren un análisis más detallado. Esta condición puede limitar la capacidad operativa de las áreas responsables, particularmente cuando deben atender varias solicitudes de manera simultánea.

	\textbf{Características de las narrativas (problema~\ref{pe:narrativas}).} Las características de las narrativas presentadas por los quejosos también influyen en el tiempo necesario para revisar las solicitudes. Debido a que los usuarios describen los acontecimientos desde su propia perspectiva, los escritos pueden variar considerablemente en extensión, estructura, claridad y nivel de detalle. Algunas narraciones incluyen información repetida, antecedentes extensos o acontecimientos que no siguen necesariamente un orden cronológico, lo que exige una lectura cuidadosa para identificar los hechos relevantes. La diversidad de estas características dificulta establecer un tiempo uniforme de revisión y aumenta la atención que debe dedicarse a determinados expedientes.

	\textbf{Falta de mecanismos para relacionar expedientes (problemas~\ref{pe:antecedentes} y~\ref{pe:duplicadas}).} La falta de mecanismos especializados para relacionar la información de diferentes expedientes puede contribuir a las dificultades para identificar antecedentes y posibles registros duplicados. Las solicitudes relacionadas con un mismo asunto no siempre presentan una redacción idéntica, debido a que los quejosos pueden describir nuevamente los acontecimientos utilizando expresiones diferentes o incorporando información adicional. Asimismo, distintos expedientes pueden compartir personas involucradas, circunstancias o características sin corresponder necesariamente al mismo hecho. Estas variaciones hacen que la identificación de relaciones entre registros requiera una comparación cuidadosa de su contenido y contexto, de modo que una queja repetida puede pasar inadvertida y registrarse como un caso nuevo.

	\textbf{Ausencia de un registro centralizado del estado de las quejas (problemas~\ref{pe:estado} y~\ref{pe:duplicadas}).} La Defensoría no cuenta con un registro único en el que se consulte la etapa en la que se encuentra cada queja, las actuaciones realizadas y los plazos que están por vencer. Sin esta información, el personal no puede identificar con facilidad los casos detenidos o próximos a vencer ni priorizar su atención, y la supervisión del conjunto de los casos depende de consultas individuales. Del lado del quejoso, la falta de un medio para consultar el avance de su queja provoca que solicite información de manera directa y, cuando considera que no ha recibido respuesta, que presente nuevamente su queja sobre los mismos hechos, lo que da origen a las duplicidades.

	\textbf{Integración limitada de la información entre áreas (problema~\ref{pe:informacion}).} La gestión de una queja involucra documentos, registros y actuaciones generados durante diferentes etapas de atención, en las que intervienen distintos integrantes del personal, cada uno con funciones y responsabilidades específicas conforme a la organización de la Defensoría \cite{ipn_manual_organizacion_ddp}. La continuidad de las actuaciones depende de que la información generada en una etapa se encuentre disponible para quienes participan posteriormente. Cuando la información se administra mediante mecanismos independientes o parcialmente integrados, resulta más complejo conservar una visión completa y actualizada del desarrollo de cada caso, y el traspaso de un expediente entre áreas requiere consultas adicionales que retrasan la atención y dificultan la coordinación de las actividades institucionales.

	\section{Propuesta}
	\label{sec:propuesta}
	Para atender los problemas identificados, se propone una plataforma web para la gestión digital de las quejas de la Defensoría que concentre en un mismo entorno el registro, la atención y el seguimiento de cada expediente, e incorpore dos componentes de inteligencia artificial como apoyo para las actividades de análisis. La plataforma no modifica el procedimiento institucional: lo representa mediante etapas y estados del expediente, de modo que cada área trabaja sobre la misma información y con los permisos que corresponden a su función.

	\begin{itemize}
		\item \textbf{Digitalización del procedimiento (problema~\ref{pe:volumen}).} Se digitaliza el procedimiento completo: el quejoso presenta su queja y sus evidencias desde un portal web y obtiene un folio de seguimiento, mientras que el personal cuenta con un panel para cada área ---Recepción, Primer Contacto, Subdefensoría y Defensora--- desde el cual recibe, valida, turna y atiende los expedientes. Las tareas repetitivas, como el registro, la generación de oficios a partir de plantillas y el envío de notificaciones por correo, se realizan desde la plataforma, lo que reduce el tiempo que el personal dedica a actividades administrativas.

		\item \textbf{Generación de resúmenes (problema~\ref{pe:narrativas}).} Se incorpora un servicio de generación de resúmenes que primero selecciona los fragmentos más relevantes de la narrativa y después los redacta en un texto breve mediante un modelo de lenguaje. El resultado se presenta al personal de Primer Contacto como un borrador que debe revisar y que no sustituye la lectura de la queja original.

		\item \textbf{Búsqueda automatizada de antecedentes (problema~\ref{pe:antecedentes}).} Se incorpora un servicio de búsqueda de antecedentes que, a partir de la narrativa de una queja, recupera los expedientes previos más parecidos combinando una búsqueda por palabras (TF-IDF) con una búsqueda por significado (MiniLM). De esta forma es posible localizar casos relacionados aunque estén redactados con palabras distintas.

		\item \textbf{Detección de posibles duplicidades (problema~\ref{pe:duplicadas}).} El mismo servicio de búsqueda señala las quejas con mayor similitud como posibles duplicidades, para que el personal determine si corresponden al mismo asunto. Además, el quejoso puede consultar el estado de su queja con su folio y su correo, lo que reduce la necesidad de presentarla de nuevo.

		\item \textbf{Seguimiento del estado de la queja (problema~\ref{pe:estado}).} Cada queja tiene un ciclo de vida con estados definidos y un historial de sus actuaciones. El quejoso consulta su avance y recibe avisos en el portal y por correo; la Subdefensoría controla los plazos de las solicitudes de información (diez días hábiles para la primera y cinco para los recordatorios); y la Defensora dispone de un panel de supervisión que señala los casos que requieren atención.

		\item \textbf{Expediente digital único (problema~\ref{pe:informacion}).} Todas las áreas trabajan sobre un expediente único almacenado en una base de datos compartida, que integra documentos, oficios, citas, notas y determinaciones. Al cambiar de etapa, el expediente se turna automáticamente al área siguiente con toda su información, y el acceso a cada operación se controla según el rol del usuario.
	\end{itemize}

	Los componentes de inteligencia artificial funcionan exclusivamente como herramientas de apoyo: sus resultados son revisados por el personal y las decisiones sobre la competencia, la investigación y la resolución de cada asunto permanecen a cargo de la Defensoría. Los objetivos, los componentes y los requerimientos de la propuesta se detallan en el Capítulo~\ref{cap:proyecto}, y su construcción en el Capítulo~\ref{cap:trabajo_realizado}.
